% Einheit 2 – Satz vom Nullprodukt – Nr. 24–31
\setcounter{aufgabe}{23}
\einheitenkopf[e2]{Einheit 2 von 4 · Satz vom Nullprodukt}
\zweigzeile{Hier lernst du, Gleichungen der Form Produkt gleich null mit dem Satz vom Nullprodukt zu lösen · neu oder schon bekannt – je nach Buch (an der Oberschule erst Kl. 10) · P10 · baut auf: lineare Gleichungen, Ausmultiplizieren und Ausklammern, Einsetzen}

\begin{aufgabe}{Ich erkenne, ob ein Produkt gleich null gesetzt ist.}
Rechne nicht. Steht links ein Produkt und rechts null? Kreuze an.
\begin{teile}
\teil $(x - 2)\cdot(x - 9) = 0$ \hfill \janein
\teil $x\cdot(x + 4) = 12$ \hfill \janein
\teil $(x + 1)\cdot(x - 6) = 0$ \hfill \janein
\teil $(x - 3)\cdot(x + 3) = 7$ \hfill \janein
\teil $x\cdot(x + 4) - 12 = 0$ \hfill \janein
\teil $x\cdot(x - 5) = 0$ \hfill \janein
\end{teile}
\end{aufgabe}

\verfahren{Gleichungen in Produktform}

\begin{aufgabe}{Ich kann eine Gleichung in Produktform mit dem Satz vom Nullprodukt lösen.}
Löse die Gleichung.
\rechenplatz{3}
\begin{gleichungsraster}[2]
\gl{(x - 2)\cdot(x - 5) = 0} & \gl{(x - 1)\cdot(x - 4) = 0} \\
\gl{(x - 6)\cdot(x - 3) = 0} & \gl{(x - 7)\cdot(x - 2) = 0} \\
\gl{(x + 4)\cdot(x - 1) = 0} & \gl{x\cdot(x + 6) = 0} \\
\gl{(x - 3)\cdot(x - 3) = 0} & \\
\end{gleichungsraster}
\end{aufgabe}

\begin{aufgabe}{Ich kann prüfen, ob eine Zahl eine Gleichung in Produktform erfüllt.}
Setze die Zahl ein und kreuze an.
\begin{teile}
\teil $x = 3$ in $(x - 3)\cdot(x + 8) = 0$ \hfill \kreuz{(wA)} \quad \kreuz{(fA)}
\teil $x = -2$ in $(x - 2)\cdot(x + 5) = 0$ \hfill \kreuz{(wA)} \quad \kreuz{(fA)}
\teil $x = -7$ in $x\cdot(x + 9) = -14$ \hfill \kreuz{(wA)} \quad \kreuz{(fA)}
\teil $x = -1$ in $x\cdot(x - 3) = -4$ \hfill \kreuz{(wA)} \quad \kreuz{(fA)}
\teil $x = -6$ in $(x + 2)\cdot(x + 3) = 12$ \hfill \kreuz{(wA)} \quad \kreuz{(fA)}
\anweisung{Welche Zahl erfüllt die Gleichung? Kreuze an.}
\teil $x\cdot(x + 7) = -10$ \hfill \kreuz{$x = 2$} \quad \kreuz{$x = 5$} \quad \kreuz{$x = -2$} \quad \kreuz{$x = -10$} \quad (P10 2025 FOR)
\end{teile}
\end{aufgabe}

\verfahren{Gleichungen, in denen jedes Glied ein $x$ hat}

\begin{aufgabe}{Ich kann eine Gleichung, in der jedes Glied ein $x$ hat, durch Ausklammern lösen.}
Löse die Gleichung.
\rechenplatz{3}
\begin{gleichungsraster}[3]
\gl{x^2 + 4x = 0} & \gl{x^2 + 5x = 0} \\
\gl{x^2 + 2x = 0} & \gl{x^2 - 6x = 0} \\
\gl{5x^2 - 15x = 0} & \\
\end{gleichungsraster}
\end{aufgabe}

\verfahren{Zwischen Produkt und Normalform wechseln}

\begin{aufgabe}{Ich kann eine Produktgleichung mit einer Zahl rechts ausmultiplizieren und ordnen.}
Multipliziere aus und bringe alles auf die linke Seite. Löse noch nicht.
\rechenplatz{2}
\begin{gleichungsraster}[2]
\gl{x\cdot(x + 3) = 10} & \gl{x\cdot(x - 4) = -3} \\
\gl{(x + 1)\cdot(x - 2) = 4} & \\
\end{gleichungsraster}
\end{aufgabe}

\begin{aufgabe}{Ich kann eine Gleichung lösen, wenn ich ihre Produktform kenne.}
Prüfe die Produktform durch Ausmultiplizieren. Gib dann die Lösungen an.
\rechenplatz{3}
\begin{teile}
\teil $x^2 + 5x + 6 = 0$ \quad Produktform: $(x + 2)\cdot(x + 3) = 0$ \hfill \feld{x_1} \quad \feld{x_2}
\teil $x^2 - 2x - 8 = 0$ \quad Produktform: $(x - 4)\cdot(x + 2) = 0$ \hfill \feld{x_1} \quad \feld{x_2}
\teil Berechne die Nullstellen der Funktion $f$ mit $f(x) = 4x^2 - 4x - 24$. Zeige dazu, dass $4x^2 - 4x - 24 = 4\cdot(x - 3)\cdot(x + 2)$ ist. (P10 2020 FOR) \hfill \feld{x_1} \quad \feld{x_2}
\end{teile}
\end{aufgabe}

\begin{aufgabe}{Ich finde den Fehler beim Satz vom Nullprodukt.}
Finde den Fehler, benenne ihn und korrigiere die Rechnung.
\begin{teile}
\teil Jana löst $x^2 = 6x$:
\rechnung{x^2 &= 6x &&\mid :x \\ x &= 6}
\teil Paul löst $(x + 5)\cdot(x - 2) = 0$:
\rechnung{x + 5 &= 0 &&\text{oder}\quad x - 2 = 0 \\ x_1 &= 5 &&\text{oder}\quad x_2 = -2}
\teil Ole löst $x\cdot(x + 1) = 12$:
\rechnung{x &= 12 &&\text{oder}\quad x + 1 = 12 \\ x_1 &= 12 &&\text{oder}\quad x_2 = 11}
\end{teile}
\end{aufgabe}

\begin{aufgabe}{Ich kann Regeln zum Satz vom Nullprodukt begründen.}
Begründe.
\begin{teile}
\teil Warum ist ein Produkt null, sobald einer der Faktoren null ist?
\teil Warum darf man bei $x^2 = 3x$ nicht einfach durch $x$ teilen?
\end{teile}
\end{aufgabe}
